\sectionheading{5.35 Input ও Output}

Program-এর কাজ সম্পন্ন করতে অনেক সময় ব্যবহারকারীর কাছ থেকে
Data গ্রহণ করতে হয় এবং Process শেষে Result প্রদর্শন করতে হয়।

\begin{center}
\textbf{Input}
\(\longrightarrow\)
\textbf{Processing}
\(\longrightarrow\)
\textbf{Output}
\end{center}

%================================================
% Input
%================================================

\booksubheading{Input কী?}

\begin{definitionbox}
Program-এ বাইরে থেকে প্রয়োজনীয় Data গ্রহণ করার প্রক্রিয়াকে
\textbf{Input} বলে।
\end{definitionbox}

Python-এ Keyboard থেকে Input নেওয়ার জন্য সাধারণত
\textenglish{\ttfamily input()}
Function ব্যবহার করা হয়।

\textbf{সাধারণ Syntax—}

\begin{codebox}
\begin{english}
\ttfamily
variable = input("Message")
\end{english}
\end{codebox}

%================================================
% Input example + behavior side-by-side
%================================================

\noindent
\begin{minipage}[t]{0.49\linewidth}
\begin{examplebox}{\textenglish{\ttfamily input()}-এর উদাহরণ}

\begin{codebox}
\begin{english}
\ttfamily
name = input("Enter your name: ")\\
print(name)
\end{english}
\end{codebox}

Input:
\begin{codebox}
\begin{english}
\ttfamily
Enter your name: Mina
\end{english}
\end{codebox}

Output:
\begin{codebox}
\begin{english}
\ttfamily
Mina
\end{english}
\end{codebox}

\end{examplebox}
\end{minipage}
\hfill
\begin{minipage}[t]{0.49\linewidth}
\begin{examplebox}{\textenglish{\ttfamily input()} কীভাবে কাজ করে}

Parentheses-এর ভিতরের Message Screen-এ দেখায় এবং Program
User-এর Input-এর জন্য অপেক্ষা করে।

\medskip

\begin{center}
\textenglish{\ttfamily input()}
\(\longrightarrow\)
\textbf{Data গ্রহণ}
\end{center}

Input করা Value Variable-এ সংরক্ষণ করা যায়।

\end{examplebox}
\end{minipage}

%================================================
% input returns str
%================================================

\booksubheading{\textenglish{\ttfamily input()}-এর Return Type}

Python-এর একটি গুরুত্বপূর্ণ নিয়ম হলো—
\textenglish{\ttfamily input()}
দিয়ে গ্রহণ করা Data সাধারণত
\textenglish{\ttfamily str}
Type হিসেবে পাওয়া যায়।

\noindent
\begin{minipage}[t]{0.49\linewidth}

\begin{codebox}
\begin{english}
\ttfamily
age = input("Enter age: ")\\
print(type(age))
\end{english}
\end{codebox}

User Input:

\begin{codebox}
\begin{english}
\ttfamily
18
\end{english}
\end{codebox}

\end{minipage}
\hfill
\begin{minipage}[t]{0.49\linewidth}

Output:

\begin{codebox}
\begin{english}
\ttfamily
<class 'str'>
\end{english}
\end{codebox}

অর্থাৎ—

\begin{center}
\textenglish{\ttfamily 18}
\(\longrightarrow\)
\textenglish{\ttfamily input()}
\(\longrightarrow\)
\textenglish{\ttfamily "18"}
\end{center}

\end{minipage}

\begin{rememberbox}
\textbf{\color{deepOrange}MCQ Trap:}
User
\textenglish{\ttfamily 18}
লিখলেও
\textenglish{\ttfamily input()}
সাধারণত সেটিকে
\textenglish{\ttfamily "18"}
অর্থাৎ String হিসেবে দেয়।
\end{rememberbox}

%================================================
% int and float side-by-side
%================================================

\booksubheading{Numeric Input নেওয়া}

Number নিয়ে Calculation করতে হলে
\textenglish{\ttfamily input()}
থেকে পাওয়া String-কে প্রয়োজনীয় Numeric Type-এ Convert করতে হয়।

\noindent
\begin{minipage}[t]{0.49\linewidth}

\begin{examplebox}{Integer Input}

\begin{codebox}
\begin{english}
\ttfamily
age = int(input("Enter age: "))
\end{english}
\end{codebox}

কার্যধারা—

\begin{center}
\textenglish{\ttfamily input()}
\(\longrightarrow\)
String
\(\longrightarrow\)
\textenglish{\ttfamily int()}
\(\longrightarrow\)
Integer
\end{center}

উদাহরণ—

\begin{codebox}
\begin{english}
\ttfamily
a = int(input("Enter number: "))\\
print(a + 5)
\end{english}
\end{codebox}

Input: \textenglish{\ttfamily 10}

Output: \textenglish{\ttfamily 15}

\end{examplebox}

\end{minipage}
\hfill
\begin{minipage}[t]{0.49\linewidth}

\begin{examplebox}{Floating-point Input}

\begin{codebox}
\begin{english}
\ttfamily
price = float(input("Enter price: "))
\end{english}
\end{codebox}

কার্যধারা—

\begin{center}
\textenglish{\ttfamily input()}
\(\longrightarrow\)
String
\(\longrightarrow\)
\textenglish{\ttfamily float()}
\(\longrightarrow\)
Float
\end{center}

উদাহরণ—

\begin{codebox}
\begin{english}
\ttfamily
x = float(input("Enter number: "))\\
print(x)
\end{english}
\end{codebox}

Input: \textenglish{\ttfamily 12.5}

Output: \textenglish{\ttfamily 12.5}

\end{examplebox}

\end{minipage}

\begin{rememberbox}

\textbf{\color{deepOrange}দ্রুত মনে রাখবে}

\begin{center}

\textbf{Text Input}
\(\longrightarrow\)
\textenglish{\ttfamily input()}
\qquad
\textbf{Integer}
\(\longrightarrow\)
\textenglish{\ttfamily int(input())}
\qquad
\textbf{Float}
\(\longrightarrow\)
\textenglish{\ttfamily float(input())}

\end{center}

\end{rememberbox}

%================================================
% Output
%================================================

\booksubheading{Output কী?}

\begin{definitionbox}
Program-এর Process করা Data, Message বা Result ব্যবহারকারীর
সামনে প্রদর্শনের প্রক্রিয়াকে \textbf{Output} বলে।
\end{definitionbox}

Python-এ Output প্রদর্শনের জন্য সাধারণত
\textenglish{\ttfamily print()}
Function ব্যবহার করা হয়।

\textbf{সাধারণ Syntax—}

\begin{codebox}
\begin{english}
\ttfamily
print(value)
\end{english}
\end{codebox}

%================================================
% Print examples 2x2 compact
%================================================

\booksubheading{\textenglish{\ttfamily print()} Function-এর ব্যবহার}

\noindent
\begin{minipage}[t]{0.49\linewidth}

\begin{examplebox}{Text Output}

\begin{codebox}
\begin{english}
\ttfamily
print("Hello")
\end{english}
\end{codebox}

Output:
\textenglish{\ttfamily Hello}

\end{examplebox}

\begin{examplebox}{Variable Output}

\begin{codebox}
\begin{english}
\ttfamily
name = "Rahim"\\
print(name)
\end{english}
\end{codebox}

Output:
\textenglish{\ttfamily Rahim}

\end{examplebox}

\end{minipage}
\hfill
\begin{minipage}[t]{0.49\linewidth}

\begin{examplebox}{Number Output}

\begin{codebox}
\begin{english}
\ttfamily
print(100)
\end{english}
\end{codebox}

Output:
\textenglish{\ttfamily 100}

\end{examplebox}

\begin{examplebox}{Expression-এর Output}

\begin{codebox}
\begin{english}
\ttfamily
print(10 + 20)
\end{english}
\end{codebox}

Output:
\textenglish{\ttfamily 30}

\end{examplebox}

\end{minipage}

%================================================
% Multiple output
%================================================

\booksubheading{Text ও Variable একসঙ্গে Output}

একটি
\textenglish{\ttfamily print()}
Function-এ Comma দিয়ে একাধিক Value দেওয়া যায়।

\begin{codebox}
\begin{english}
\ttfamily
name = "Mina"\\
age = 18\\
print("Name:", name)\\
print("Age:", age)
\end{english}
\end{codebox}

\textbf{Output:}
\quad
\textenglish{\ttfamily Name: Mina}
\qquad
\textenglish{\ttfamily Age: 18}

%================================================
% Full program
%================================================

\booksubheading{একটি সম্পূর্ণ Input--Output Program}

\begin{examplebox}{নাম ও বয়স গ্রহণ}

\begin{codebox}
\begin{english}
\ttfamily
name = input("Enter your name: ")\\
age = int(input("Enter your age: "))\\
print("Name:", name)\\
print("Age:", age)
\end{english}
\end{codebox}

\textbf{Input:}
\quad
\textenglish{\ttfamily Rahim, 18}

\qquad
\textbf{Output:}
\quad
\textenglish{\ttfamily Name: Rahim, Age: 18}

\end{examplebox}

\begin{center}
\textbf{Input}
\(\longrightarrow\)
\textbf{Variable}
\(\longrightarrow\)
\textbf{Processing}
\(\longrightarrow\)
\textenglish{\ttfamily print()}
\(\longrightarrow\)
\textbf{Output}
\end{center}

%================================================
% Addition + trap side by side
%================================================

\booksubheading{Numeric Input-এ Conversion কেন গুরুত্বপূর্ণ?}

\noindent
\begin{minipage}[t]{0.49\linewidth}

\begin{examplebox}{সঠিক: Integer Conversion}

\begin{codebox}
\begin{english}
\ttfamily
a = int(input())\\
b = int(input())\\
print(a + b)
\end{english}
\end{codebox}

Input:
\textenglish{\ttfamily 10, 20}

\medskip

Output:
\textenglish{\ttfamily 30}

\end{examplebox}

\end{minipage}
\hfill
\begin{minipage}[t]{0.49\linewidth}

\begin{examplebox}{Conversion না করলে}

\begin{codebox}
\begin{english}
\ttfamily
a = input()\\
b = input()\\
print(a + b)
\end{english}
\end{codebox}

Input:
\textenglish{\ttfamily 10, 20}

\medskip

Output:
\textenglish{\ttfamily 1020}

\end{examplebox}

\end{minipage}

\begin{rememberbox}

\begin{center}

\textenglish{\ttfamily "10" + "20"}
\(\longrightarrow\)
\textenglish{\ttfamily "1020"}
\qquad
কিন্তু
\qquad
\textenglish{\ttfamily 10 + 20}
\(\longrightarrow\)
\textenglish{\ttfamily 30}

\end{center}

\end{rememberbox}

%================================================
% Complete addition
%================================================

\booksubheading{দুটি সংখ্যা Input নিয়ে যোগ}

\begin{codebox}
\begin{english}
\ttfamily
a = int(input("Enter first number: "))\\
b = int(input("Enter second number: "))\\
total = a + b\\
print("Sum =", total)
\end{english}
\end{codebox}

Input:
\textenglish{\ttfamily 10, 20}
\qquad
Output:
\textenglish{\ttfamily Sum = 30}

%================================================
% sep and end side by side
%================================================

\booksubheading{\textenglish{\ttfamily sep} ও \textenglish{\ttfamily end} Parameter}

\noindent
\begin{minipage}[t]{0.49\linewidth}

\begin{examplebox}{\textenglish{\ttfamily sep}}

একাধিক Value-এর \textbf{মাঝখানে} কী থাকবে তা নির্ধারণ করে।

\begin{codebox}
\begin{english}
\ttfamily
print("ICT", "Python", "HSC", sep="-")
\end{english}
\end{codebox}

Output:
\begin{codebox}
\begin{english}
\ttfamily
ICT-Python-HSC
\end{english}
\end{codebox}

Default
\textenglish{\ttfamily sep}
হলো Space।

\end{examplebox}

\end{minipage}
\hfill
\begin{minipage}[t]{0.49\linewidth}

\begin{examplebox}{\textenglish{\ttfamily end}}

\textenglish{\ttfamily print()}
-এর \textbf{শেষে} কী থাকবে তা নির্ধারণ করে।

\begin{codebox}
\begin{english}
\ttfamily
print("Hello", end=" ")\\
print("Python")
\end{english}
\end{codebox}

Output:
\begin{codebox}
\begin{english}
\ttfamily
Hello Python
\end{english}
\end{codebox}

Default
\textenglish{\ttfamily end}
হলো Newline।

\end{examplebox}

\end{minipage}

\begin{tabularx}
{\linewidth}
{
|>{\bfseries\RaggedRight\arraybackslash}p{0.23\linewidth}
|Y
|Y|
}
\hline
\rowcolor{softBlue}
\multicolumn{1}{|c|}{\textbf{বিষয়}}
&
\multicolumn{1}{c|}{\textenglish{\ttfamily sep}}
&
\multicolumn{1}{c|}{\textenglish{\ttfamily end}}
\\
\hline

কাজ
&
Valueগুলোর মাঝখানে কী থাকবে
&
\textenglish{\ttfamily print()}
শেষে কী থাকবে
\\
\hline

Default
&
Space
&
Newline
\\
\hline

উদাহরণ
&
\textenglish{\ttfamily sep="-"}
&
\textenglish{\ttfamily end=" "}
\\
\hline
\end{tabularx}

%================================================
% Escape sequences side-by-side
%================================================

\booksubheading{Escape Sequence}

Output-এর বিশেষ Formatting-এর জন্য Backslash দিয়ে Escape Sequence
ব্যবহার করা যায়।

\noindent
\begin{minipage}[t]{0.49\linewidth}

\begin{examplebox}{New Line — \textenglish{\ttfamily \textbackslash n}}

\begin{codebox}
\begin{english}
\ttfamily
print("ICT\textbackslash nPython")
\end{english}
\end{codebox}

Output:

\begin{codebox}
\begin{english}
\ttfamily
ICT\\
Python
\end{english}
\end{codebox}

\end{examplebox}

\end{minipage}
\hfill
\begin{minipage}[t]{0.49\linewidth}

\begin{examplebox}{Tab — \textenglish{\ttfamily \textbackslash t}}

\begin{codebox}
\begin{english}
\ttfamily
print("Name\textbackslash tMarks")
\end{english}
\end{codebox}

Output-এর ধারণা:

\begin{codebox}
\begin{english}
\ttfamily
Name    Marks
\end{english}
\end{codebox}

\end{examplebox}

\end{minipage}

\begin{tabularx}
{\linewidth}
{|>{\bfseries}p{0.30\linewidth}|Y|}
\hline
\rowcolor{softBlue}
\multicolumn{1}{|c|}{\textbf{Escape Sequence}}
&
\multicolumn{1}{c|}{\textbf{কাজ}}
\\
\hline

\textenglish{\ttfamily \textbackslash n}
&
নতুন Line তৈরি করে।
\\
\hline

\textenglish{\ttfamily \textbackslash t}
&
Tab Space তৈরি করে।
\\
\hline
\end{tabularx}

%================================================
% f-string
%================================================

\booksubheading{Formatted Output — f-string}

Variable-এর Value Text-এর সঙ্গে সুন্দরভাবে প্রদর্শনের জন্য
Python-এর \textbf{f-string} ব্যবহার করা যায়।

\begin{codebox}
\begin{english}
\ttfamily
name = "Rahim"\\
age = 18\\
print(f"My name is \{name\} and I am \{age\} years old.")
\end{english}
\end{codebox}

\textbf{Output:}
\quad
\textenglish{\ttfamily My name is Rahim and I am 18 years old.}

\begin{rememberbox}
\textenglish{\ttfamily f"..."}
-এর ভিতরে
\textenglish{\ttfamily \{variable\}}
ব্যবহার করলে Variable-এর Value Output-এ বসে যায়।
\end{rememberbox}

%================================================
% Input vs Output
%================================================

\booksubheading{Input ও Output-এর পার্থক্য}

\begin{tabularx}
{\linewidth}
{
|>{\bfseries\RaggedRight\arraybackslash}p{0.23\linewidth}
|Y
|Y|
}
\hline
\rowcolor{softBlue}

\multicolumn{1}{|c|}{\textbf{বিষয়}}
&
\multicolumn{1}{c|}{\textbf{Input}}
&
\multicolumn{1}{c|}{\textbf{Output}}
\\
\hline

উদ্দেশ্য
&
Data গ্রহণ
&
Result প্রদর্শন
\\
\hline

প্রধান Function
&
\textenglish{\ttfamily input()}
&
\textenglish{\ttfamily print()}
\\
\hline

Data Flow
&
User \(\longrightarrow\) Program
&
Program \(\longrightarrow\) User
\\
\hline

উদাহরণ
&
নাম, বয়স, সংখ্যা গ্রহণ
&
Message বা Result দেখানো
\\
\hline
\end{tabularx}

%================================================
% Important functions
%================================================

\booksubheading{গুরুত্বপূর্ণ Function এক নজরে}

\begin{tabularx}
{\linewidth}
{|>{\bfseries\RaggedRight\arraybackslash}p{0.30\linewidth}|Y|}
\hline
\rowcolor{softBlue}

\multicolumn{1}{|c|}{\textbf{Function}}
&
\multicolumn{1}{c|}{\textbf{কাজ}}
\\
\hline

\textenglish{\ttfamily input()}
&
User-এর কাছ থেকে Input গ্রহণ
\\
\hline

\textenglish{\ttfamily print()}
&
Output প্রদর্শন
\\
\hline

\textenglish{\ttfamily int()}
&
Integer-এ Convert
\\
\hline

\textenglish{\ttfamily float()}
&
Float-এ Convert
\\
\hline

\textenglish{\ttfamily type()}
&
Data Type নির্ণয়
\\
\hline
\end{tabularx}

%================================================
% Common mistakes
%================================================

\booksubheading{সাধারণ ভুল}

\begin{compacttwocolpoints}

\item
\textenglish{\ttfamily input()}
থেকে পাওয়া Numeric-looking Value-কে সরাসরি Number মনে করা।

\item Calculation-এর আগে প্রয়োজনীয়
\textenglish{\ttfamily int()}
বা
\textenglish{\ttfamily float()}
Conversion না করা।

\item
\textenglish{\ttfamily print("a")}
ও
\textenglish{\ttfamily print(a)}
-কে একই মনে করা।

\item
\textenglish{\ttfamily sep}
ও
\textenglish{\ttfamily end}
-এর কাজ গুলিয়ে ফেলা।

\item
\textenglish{\ttfamily \textbackslash n}
ও
\textenglish{\ttfamily \textbackslash t}
-এর কাজ গুলিয়ে ফেলা।

\item Input Prompt-এর Message-কে User-এর Input Value-এর অংশ মনে করা।

\item Text লেখার Quotation Mark ভুলে যাওয়া।

\end{compacttwocolpoints}

%================================================
% Final Reminder
%================================================

\begin{rememberbox}

\textbf{\color{deepOrange}মনে রাখবে}

\begin{itemize}

\item
\textenglish{\ttfamily input()}
\(\longrightarrow\)
User-এর কাছ থেকে Data গ্রহণ।

\item
\textenglish{\ttfamily input()}
-এর Result সাধারণত
\textenglish{\ttfamily str}।

\item
Integer Input
\(\longrightarrow\)
\textenglish{\ttfamily int(input())}

\item
Float Input
\(\longrightarrow\)
\textenglish{\ttfamily float(input())}

\item
\textenglish{\ttfamily print()}
\(\longrightarrow\)
Output প্রদর্শন।

\item
Comma দিয়ে
\textenglish{\ttfamily print()}
-এ একাধিক Value দেওয়া যায়।

\item
\textenglish{\ttfamily sep}
\(\longrightarrow\)
Value-এর মাঝে কী থাকবে।

\item
\textenglish{\ttfamily end}
\(\longrightarrow\)
Output-এর শেষে কী থাকবে।

\item
\textenglish{\ttfamily \textbackslash n}
\(\longrightarrow\)
New Line।

\item
\textenglish{\ttfamily \textbackslash t}
\(\longrightarrow\)
Tab Space।

\item
f-string
\(\longrightarrow\)
Formatted Output।

\end{itemize}

\begin{center}

\textbf{Input}
\(\longrightarrow\)
\textenglish{\ttfamily input()}
\qquad
\textbf{Output}
\(\longrightarrow\)
\textenglish{\ttfamily print()}

\medskip

\textenglish{\ttfamily sep}
\(\longrightarrow\)
\textbf{মাঝে}
\qquad
\textenglish{\ttfamily end}
\(\longrightarrow\)
\textbf{শেষে}

\end{center}

\end{rememberbox}

\small

%================================================
% 5.36 Python Operator
%================================================

